\IfFileExists{stacks-project.cls}{%
\documentclass{stacks-project}
}{%
\documentclass{amsart}
}

% For dealing with references we use the comment environment
\usepackage{verbatim}
\newenvironment{reference}{\comment}{\endcomment}
%\newenvironment{reference}{}{}
\newenvironment{slogan}{\comment}{\endcomment}
\newenvironment{history}{\comment}{\endcomment}

% For commutative diagrams we use Xy-pic
\usepackage[all]{xy}

% We use 2cell for 2-commutative diagrams.
\xyoption{2cell}
\UseAllTwocells

% We use multicol for the list of chapters between chapters
\usepackage{multicol}

% This is generally recommended for better output
\usepackage{lmodern}
\usepackage[T1]{fontenc}

% For cross-file-references
\usepackage{xr-hyper}

% Package for hypertext links:
\usepackage{hyperref}

% For any local file, say "hello.tex" you want to link to please
% use \externaldocument[hello-]{hello}
\externaldocument[introduction-]{introduction}
\externaldocument[conventions-]{conventions}
\externaldocument[sets-]{sets}
\externaldocument[categories-]{categories}
\externaldocument[topology-]{topology}
\externaldocument[sheaves-]{sheaves}
\externaldocument[sites-]{sites}
\externaldocument[stacks-]{stacks}
\externaldocument[fields-]{fields}
\externaldocument[algebra-]{algebra}
\externaldocument[brauer-]{brauer}
\externaldocument[homology-]{homology}
\externaldocument[derived-]{derived}
\externaldocument[simplicial-]{simplicial}
\externaldocument[more-algebra-]{more-algebra}
\externaldocument[smoothing-]{smoothing}
\externaldocument[modules-]{modules}
\externaldocument[sites-modules-]{sites-modules}
\externaldocument[injectives-]{injectives}
\externaldocument[cohomology-]{cohomology}
\externaldocument[sites-cohomology-]{sites-cohomology}
\externaldocument[dga-]{dga}
\externaldocument[dpa-]{dpa}
\externaldocument[sdga-]{sdga}
\externaldocument[hypercovering-]{hypercovering}
\externaldocument[schemes-]{schemes}
\externaldocument[constructions-]{constructions}
\externaldocument[properties-]{properties}
\externaldocument[morphisms-]{morphisms}
\externaldocument[coherent-]{coherent}
\externaldocument[divisors-]{divisors}
\externaldocument[limits-]{limits}
\externaldocument[varieties-]{varieties}
\externaldocument[topologies-]{topologies}
\externaldocument[descent-]{descent}
\externaldocument[perfect-]{perfect}
\externaldocument[more-morphisms-]{more-morphisms}
\externaldocument[flat-]{flat}
\externaldocument[groupoids-]{groupoids}
\externaldocument[more-groupoids-]{more-groupoids}
\externaldocument[etale-]{etale}
\externaldocument[chow-]{chow}
\externaldocument[intersection-]{intersection}
\externaldocument[pic-]{pic}
\externaldocument[weil-]{weil}
\externaldocument[adequate-]{adequate}
\externaldocument[dualizing-]{dualizing}
\externaldocument[duality-]{duality}
\externaldocument[discriminant-]{discriminant}
\externaldocument[derham-]{derham}
\externaldocument[local-cohomology-]{local-cohomology}
\externaldocument[algebraization-]{algebraization}
\externaldocument[curves-]{curves}
\externaldocument[resolve-]{resolve}
\externaldocument[models-]{models}
\externaldocument[functors-]{functors}
\externaldocument[equiv-]{equiv}
\externaldocument[pione-]{pione}
\externaldocument[etale-cohomology-]{etale-cohomology}
\externaldocument[proetale-]{proetale}
\externaldocument[relative-cycles-]{relative-cycles}
\externaldocument[more-etale-]{more-etale}
\externaldocument[trace-]{trace}
\externaldocument[crystalline-]{crystalline}
\externaldocument[spaces-]{spaces}
\externaldocument[spaces-properties-]{spaces-properties}
\externaldocument[spaces-morphisms-]{spaces-morphisms}
\externaldocument[decent-spaces-]{decent-spaces}
\externaldocument[spaces-cohomology-]{spaces-cohomology}
\externaldocument[spaces-limits-]{spaces-limits}
\externaldocument[spaces-divisors-]{spaces-divisors}
\externaldocument[spaces-over-fields-]{spaces-over-fields}
\externaldocument[spaces-topologies-]{spaces-topologies}
\externaldocument[spaces-descent-]{spaces-descent}
\externaldocument[spaces-perfect-]{spaces-perfect}
\externaldocument[spaces-more-morphisms-]{spaces-more-morphisms}
\externaldocument[spaces-flat-]{spaces-flat}
\externaldocument[spaces-groupoids-]{spaces-groupoids}
\externaldocument[spaces-more-groupoids-]{spaces-more-groupoids}
\externaldocument[bootstrap-]{bootstrap}
\externaldocument[spaces-pushouts-]{spaces-pushouts}
\externaldocument[spaces-chow-]{spaces-chow}
\externaldocument[groupoids-quotients-]{groupoids-quotients}
\externaldocument[spaces-more-cohomology-]{spaces-more-cohomology}
\externaldocument[spaces-simplicial-]{spaces-simplicial}
\externaldocument[spaces-duality-]{spaces-duality}
\externaldocument[formal-spaces-]{formal-spaces}
\externaldocument[restricted-]{restricted}
\externaldocument[spaces-resolve-]{spaces-resolve}
\externaldocument[formal-defos-]{formal-defos}
\externaldocument[defos-]{defos}
\externaldocument[cotangent-]{cotangent}
\externaldocument[examples-defos-]{examples-defos}
\externaldocument[algebraic-]{algebraic}
\externaldocument[examples-stacks-]{examples-stacks}
\externaldocument[stacks-sheaves-]{stacks-sheaves}
\externaldocument[criteria-]{criteria}
\externaldocument[artin-]{artin}
\externaldocument[quot-]{quot}
\externaldocument[stacks-properties-]{stacks-properties}
\externaldocument[stacks-morphisms-]{stacks-morphisms}
\externaldocument[stacks-limits-]{stacks-limits}
\externaldocument[stacks-cohomology-]{stacks-cohomology}
\externaldocument[stacks-perfect-]{stacks-perfect}
\externaldocument[stacks-introduction-]{stacks-introduction}
\externaldocument[stacks-more-morphisms-]{stacks-more-morphisms}
\externaldocument[stacks-geometry-]{stacks-geometry}
\externaldocument[moduli-]{moduli}
\externaldocument[moduli-curves-]{moduli-curves}
\externaldocument[examples-]{examples}
\externaldocument[exercises-]{exercises}
\externaldocument[guide-]{guide}
\externaldocument[desirables-]{desirables}
\externaldocument[coding-]{coding}
\externaldocument[obsolete-]{obsolete}
\externaldocument[fdl-]{fdl}
\externaldocument[index-]{index}

% Theorem environments.
%
\theoremstyle{plain}
\newtheorem{theorem}[subsection]{Theorem}
\newtheorem{proposition}[subsection]{Proposition}
\newtheorem{lemma}[subsection]{Lemma}

\theoremstyle{definition}
\newtheorem{definition}[subsection]{Definition}
\newtheorem{example}[subsection]{Example}
\newtheorem{exercise}[subsection]{Exercise}
\newtheorem{situation}[subsection]{Situation}

\theoremstyle{remark}
\newtheorem{remark}[subsection]{Remark}
\newtheorem{remarks}[subsection]{Remarks}

\numberwithin{equation}{subsection}

% Macros
%
\def\lim{\mathop{\mathrm{lim}}\nolimits}
\def\colim{\mathop{\mathrm{colim}}\nolimits}
\def\Spec{\mathop{\mathrm{Spec}}}
\def\Hom{\mathop{\mathrm{Hom}}\nolimits}
\def\Ext{\mathop{\mathrm{Ext}}\nolimits}
\def\SheafHom{\mathop{\mathcal{H}\!\mathit{om}}\nolimits}
\def\SheafExt{\mathop{\mathcal{E}\!\mathit{xt}}\nolimits}
\def\Sch{\mathit{Sch}}
\def\Mor{\mathop{\mathrm{Mor}}\nolimits}
\def\Ob{\mathop{\mathrm{Ob}}\nolimits}
\def\Sh{\mathop{\mathit{Sh}}\nolimits}
\def\NL{\mathop{N\!L}\nolimits}
\def\CH{\mathop{\mathrm{CH}}\nolimits}
\def\proetale{{pro\text{-}\acute{e}tale}}
\def\etale{{\acute{e}tale}}
\def\QCoh{\mathit{QCoh}}
\def\Ker{\mathop{\mathrm{Ker}}}
\def\Im{\mathop{\mathrm{Im}}}
\def\Coker{\mathop{\mathrm{Coker}}}
\def\Coim{\mathop{\mathrm{Coim}}}

% Boxtimes
%
\DeclareMathSymbol{\boxtimes}{\mathbin}{AMSa}{"02}

%
% Macros for moduli stacks/spaces
%
\def\QCohstack{\mathcal{QC}\!\mathit{oh}}
\def\Cohstack{\mathcal{C}\!\mathit{oh}}
\def\Spacesstack{\mathcal{S}\!\mathit{paces}}
\def\Quotfunctor{\mathrm{Quot}}
\def\Hilbfunctor{\mathrm{Hilb}}
\def\Curvesstack{\mathcal{C}\!\mathit{urves}}
\def\Polarizedstack{\mathcal{P}\!\mathit{olarized}}
\def\Complexesstack{\mathcal{C}\!\mathit{omplexes}}
% \Pic is the operator that assigns to X its picard group, usage \Pic(X)
% \Picardstack_{X/B} denotes the Picard stack of X over B
% \Picardfunctor_{X/B} denotes the Picard functor of X over B
\def\Pic{\mathop{\mathrm{Pic}}\nolimits}
\def\Picardstack{\mathcal{P}\!\mathit{ic}}
\def\Picardfunctor{\mathrm{Pic}}
\def\Deformationcategory{\mathcal{D}\!\mathit{ef}}
\usepackage{cleveref}
\crefname{originaldefinition}{original definition}{original definitions}
\Crefname{originaldefinition}{Original definition}{Original definitions}
\makeatletter
\newlabel{situation-q}{{082M}{source}{}{sourcepointer-1}{}}
\newlabel{situation-q@cref}{{[originaldefinition][0][]082M}{[1][0][]source}{}{}{}}
\makeatother
\crefname{originallemma}{original lemma}{original lemmas}
\Crefname{originallemma}{Original lemma}{Original lemmas}
\makeatletter
\newlabel{spaces-limits-lemma-directed-inverse-system-has-limit}{{07SF}{source}{}{sourcepointer-2}{}}
\newlabel{spaces-limits-lemma-directed-inverse-system-has-limit@cref}{{[originallemma][0][]07SF}{[1][0][]source}{}{}{}}
\makeatother
\crefname{originallemma}{original lemma}{original lemmas}
\Crefname{originallemma}{Original lemma}{Original lemmas}
\makeatletter
\newlabel{spaces-limits-lemma-descend-modules-finite-presentation}{{07V7}{source}{}{sourcepointer-3}{}}
\newlabel{spaces-limits-lemma-descend-modules-finite-presentation@cref}{{[originallemma][0][]07V7}{[1][0][]source}{}{}{}}
\makeatother
\crefname{originallemma}{original lemma}{original lemmas}
\Crefname{originallemma}{Original lemma}{Original lemmas}
\makeatletter
\newlabel{spaces-properties-lemma-quasi-compact-affine-cover}{{03H6}{source}{}{sourcepointer-4}{}}
\newlabel{spaces-properties-lemma-quasi-compact-affine-cover@cref}{{[originallemma][0][]03H6}{[1][0][]source}{}{}{}}
\makeatother
\crefname{originallemma}{original lemma}{original lemmas}
\Crefname{originallemma}{Original lemma}{Original lemmas}
\makeatletter
\newlabel{algebra-lemma-module-map-property-in-colimit}{{05LI}{source}{}{sourcepointer-5}{}}
\newlabel{algebra-lemma-module-map-property-in-colimit@cref}{{[originallemma][0][]05LI}{[1][0][]source}{}{}{}}
\makeatother
\crefname{originallemma}{original lemma}{original lemmas}
\Crefname{originallemma}{Original lemma}{Original lemmas}
\makeatletter
\newlabel{algebra-lemma-flat-finite-presentation-limit-flat}{{02JO}{source}{}{sourcepointer-6}{}}
\newlabel{algebra-lemma-flat-finite-presentation-limit-flat@cref}{{[originallemma][0][]02JO}{[1][0][]source}{}{}{}}
\makeatother
\title{The finite-presentation quotient functor is limit preserving}
\author{The Stacks Project Authors}
\date{stacks-project:quot.tex:local-snapshot-20261012; GNU FDL 1.2 or later}
\begin{document}
\maketitle
\setcounter{subsection}{1}
\noindent\textit{Editorial note.} All delivered mathematical text is the author's own, copied from the pinned Stacks Project chapter file quot.tex; no statement or proof step was rewritten, reordered or removed. The counted span is the author environment and its own proof environment. References inside the span are handed to the assembler as navigation stubs that carry the canonical tag of the cited tag only; whether each of them is load-bearing is recorded on the stub itself.\par\medskip
\noindent\textit{Allowed standard prerequisites (editorial).}\begin{itemize}
\item Situation situation-q (canonical tag 082M). cited at statement line 2164. Standing situation: f : X -> B with the finite presentation quotient functor Q\textasciicircum{}fp\_\{F/X/B\}.
\item Lemma spaces-limits-lemma-directed-inverse-system-has-limit (canonical tag 07SF). cited at proof line 2181. X\_T = lim X\_\{T\_i\} for the directed limit of affine schemes.
\item Lemma spaces-limits-lemma-descend-modules-finite-presentation (canonical tag 07V7). cited at proof line 2185. The quotient map F\_T -> Q descends to F\_\{T\_i\} -> Q\_i of finite presentation.
\item Lemma spaces-properties-lemma-quasi-compact-affine-cover (canonical tag 03H6). cited at proof line 2195. Affine etale cover U -> X reduces surjectivity and flatness checks to the affine case.
\item Lemma algebra-lemma-module-map-property-in-colimit (canonical tag 05LI). cited at proof line 2208. Surjectivity at the colimit holds already at a finite stage.
\item Lemma algebra-lemma-flat-finite-presentation-limit-flat (canonical tag 02JO). cited at proof line 2210. Flatness at the colimit holds already at a finite stage.
\end{itemize}\medskip
\ifdefined\hypertarget\hypertarget{sourcepointer-1}{}\fi
\noindent\textit{Source reference.} Original definition 082M: sources/stacks-chow/tags.txt, source lines 10257--10257. Cited at statement line 2164: the lemma is phrased in this situation; load-bearing.\par\medskip
\ifdefined\hypertarget\hypertarget{sourcepointer-2}{}\fi
\noindent\textit{Source reference.} Original lemma 07SF: sources/stacks-chow/tags.txt, source lines 9900--9900. Cited at proof line 2181: the base change X\_T is the directed limit of the X\_\{T\_i\}; load-bearing.\par\medskip
\ifdefined\hypertarget\hypertarget{sourcepointer-3}{}\fi
\noindent\textit{Source reference.} Original lemma 07V7: sources/stacks-chow/tags.txt, source lines 9997--9997. Cited at proof line 2185: the quotient map F\_T -> Q descends to some F\_\{T\_i\} -> Q\_i of finite presentation; load-bearing.\par\medskip
\ifdefined\hypertarget\hypertarget{sourcepointer-4}{}\fi
\noindent\textit{Source reference.} Original lemma 03H6: sources/stacks-chow/tags.txt, source lines 4555--4555. Cited at proof line 2195: choose the affine etale cover U -> X to reduce to the affine case; load-bearing.\par\medskip
\ifdefined\hypertarget\hypertarget{sourcepointer-5}{}\fi
\noindent\textit{Source reference.} Original lemma 05LI: sources/stacks-chow/tags.txt, source lines 7237--7237. Cited at proof line 2208: surjectivity of M\_i -> N\_i holds after enlarging i; load-bearing.\par\medskip
\ifdefined\hypertarget\hypertarget{sourcepointer-6}{}\fi
\noindent\textit{Source reference.} Original lemma 02JO: sources/stacks-chow/tags.txt, source lines 3343--3343. Cited at proof line 2210: flatness of N over A holds after enlarging i; load-bearing.\par\medskip
\begin{lemma}
\label{lemma-q-limit-preserving}
In Situation \ref{situation-q} assume also that
(a) $f$ is quasi-compact and quasi-separated and
(b) $\mathcal{F}$ is of finite presentation.
Then the functor $\text{Q}^{fp}_{\mathcal{F}/X/B}$
is limit preserving in the following sense: If $T = \lim T_i$ is a
directed limit of affine schemes over $B$, then
$\text{Q}^{fp}_{\mathcal{F}/X/B}(T) =
\colim \text{Q}^{fp}_{\mathcal{F}/X/B}(T_i)$.
\end{lemma}

\begin{proof}
Let $T = \lim T_i$ be as in the statement of the lemma.
Choose $i_0 \in I$ and replace $I$ by $\{i \in I \mid i \geq i_0\}$.
We may set $B = S = T_{i_0}$ and we may replace $X$ by $X_{T_0}$
and $\mathcal{F}$ by the pullback to $X_{T_0}$. Then
$X_T = \lim X_{T_i}$, see
Limits of Spaces, Lemma
\ref{spaces-limits-lemma-directed-inverse-system-has-limit}.
Let $\mathcal{F}_T \to \mathcal{Q}$ be an element of
$\text{Q}^{fp}_{\mathcal{F}/X/B}(T)$. By
Limits of Spaces, Lemma
\ref{spaces-limits-lemma-descend-modules-finite-presentation}
there exists an $i$ and a map $\mathcal{F}_{T_i} \to \mathcal{Q}_i$
of $\mathcal{O}_{X_{T_i}}$-modules of finite presentation whose
pullback to $X_T$ is the given quotient map.

\medskip\noindent
We still have to check that, after possibly increasing $i$, the map
$\mathcal{F}_{T_i} \to \mathcal{Q}_i$ is surjective and $\mathcal{Q}_i$
is flat over $T_i$. To do this, choose an affine scheme $U$ and a
surjective \'etale morphism $U \to X$ (see Properties of Spaces,
Lemma \ref{spaces-properties-lemma-quasi-compact-affine-cover}).
We may check surjectivity and flatness over $T_i$ after pulling
back to the \'etale cover $U_{T_i} \to X_{T_i}$ (by definition).
This reduces us to the case where $X = \Spec(B_0)$ is an affine scheme of
finite presentation over $B = S = T_0 = \Spec(A_0)$.
Writing $T_i = \Spec(A_i)$, then $T = \Spec(A)$ with $A = \colim A_i$
we have reached the following algebra problem. Let $M_i \to N_i$
be a map of finitely presented $B_0 \otimes_{A_0} A_i$-modules
such that $M_i \otimes_{A_i} A \to N_i \otimes_{A_i} A$ is surjective
and $N_i \otimes_{A_i} A$ is flat over $A$. Show that for some $i' \geq i$
$M_i \otimes_{A_i} A_{i'} \to N_i \otimes_{A_i} A_{i'}$ is surjective
and $N_i \otimes_{A_i} A_{i'}$ is flat over $A$.
The first follows from
Algebra, Lemma \ref{algebra-lemma-module-map-property-in-colimit}
and the second from
Algebra, Lemma \ref{algebra-lemma-flat-finite-presentation-limit-flat}.
\end{proof}

\end{document}
